\documentclass[11pt,a4paper]{article}

% ---------------------------------------------------------
% PREAMBLE
% ---------------------------------------------------------
\usepackage[utf8]{inputenc}
\usepackage[T1]{fontenc}
\usepackage{lmodern}
\usepackage{amsmath,amssymb,amsfonts,bm}
\usepackage{geometry}
\usepackage{graphicx}
\usepackage{hyperref}
\usepackage{cite}
\usepackage{authblk}
\usepackage{physics}
\usepackage{mathtools}
\usepackage{enumitem}

\geometry{margin=2.4cm}

\title{\textbf{Companion L (v49.1)}\\[0.2cm]
\textbf{The Growth Index in the Relaxation-Driven Cyclic Cosmology}\\[0.2cm]
\large Foundations of Relaxation-Driven Cyclic Cosmology (RDCC)}

\author[1]{Michael Lehmann}
\affil[1]{\small Independent Researcher, Relaxation-Driven Cyclic Cosmology (RDCC)\\
\texttt{mi.lehmann@gmx.de}}

\date{September 2026}

\begin{document}
\maketitle

\begin{abstract}
The growth index \(\gamma\) provides a compact description of the dynamical
evolution of cosmic structure. In the standard cosmological model, \(\gamma\)
is nearly constant, reflecting the scale-independent growth of density
perturbations. In previous versions of RDCC the growth index was written with
an inserted relaxation scale. In the present Phase-1 integrated version this
scale is inherited from the freeze-out geometry of Companion XII,
\(k_{\mathrm{rel}}=a_{\mathrm{fr}}H(a_{\mathrm{fr}})\). The resulting
scale-dependent and time-dependent growth index traces the same geometric
transition that governs tensor damping, gravitational slip, and the Boltzmann
hierarchy. This Companion develops the continuous description of \(\gamma\) in
RDCC and identifies its signatures in redshift-space distortions, weak lensing,
and velocity fields.
\end{abstract}

\tableofcontents
\newpage

% ---------------------------------------------------------
\section{Introduction}
% ---------------------------------------------------------

The growth index \(\gamma\) is defined through the relation
\begin{equation}
  f(a)
  =
  \Omega_{m}(a)^{\gamma},
\end{equation}
where \(f(a)\) is the linear growth rate. In the standard cosmological model
the growth index is nearly constant, with \(\gamma\approx 0.55\) for a Universe
dominated by cold dark matter and a cosmological constant.

In earlier versions of RDCC the gravitational potential was modified by a
phenomenological scale \(k_{\mathrm{rel}}\). That description remains useful as
a fitting template. In v49.1 the same scale is no longer independent: it is
the comoving horizon scale at freeze-out. The growth index therefore becomes a
probe of the freeze-out structure of the RDCC gravitational field rather than
an extra parameter.

% ---------------------------------------------------------
\section{Connection to the Freeze-Out Geometry}
% ---------------------------------------------------------

Companion XII established the relaxation scale
\begin{equation}
  k_{\mathrm{rel}}
  =
  a_{\mathrm{fr}} H(a_{\mathrm{fr}}),
\end{equation}
where
\begin{equation}
  a_{\mathrm{fr}}
  \simeq
  10^{-4}
\end{equation}
marks the freeze-out epoch of the global relaxation field.

The growth hierarchy inherits this scale from the same relaxation geometry that
governs tensor damping and gravitational slip. The growth index therefore
probes the freeze-out structure of the RDCC Universe rather than an independent
phenomenological scale.

The resulting causal chain is
\begin{equation}
  \alpha
  \;\rightarrow\;
  a_{\mathrm{fr}}
  \;\rightarrow\;
  k_{\mathrm{rel}}
  \;\rightarrow\;
  D(k,a)
  \;\rightarrow\;
  f(k,a)
  \;\rightarrow\;
  \gamma(k,a).
\end{equation}

% ---------------------------------------------------------
\section{The RDCC Growth Rate}
% ---------------------------------------------------------

The linear growth rate is defined as
\begin{equation}
  f(k,a)
  =
  \frac{d\ln D(k,a)}{d\ln a},
\end{equation}
where \(D(k,a)\) is the growth factor.

In RDCC the growth factor evolves as
\begin{equation}
  D_{\mathrm{RDCC}}(k,a)
  =
  D(a)
  \left[
    1
    -
    \chi_{D}
    \left(
      \frac{k}{k_{\mathrm{rel}}}
    \right)^{\alpha}
  \right].
\end{equation}
The relaxation scale appearing in the growth factor is not introduced as an
independent parameter. Within the Phase-1 integrated RDCC framework it is
determined by the freeze-out geometry according to
\begin{equation}
  k_{\mathrm{rel}}
  =
  a_{\mathrm{fr}} H(a_{\mathrm{fr}}).
\end{equation}
Consequently, the onset of growth suppression is linked directly to the
physical freeze-out of the relaxation sector.

The growth rate becomes
\begin{equation}
  f_{\mathrm{RDCC}}(k,a)
  =
  f(a)
  \left[
    1
    -
    \chi_{f}
    \left(
      \frac{k}{k_{\mathrm{rel}}}
    \right)^{\alpha}
  \right].
\end{equation}
The coefficients
\begin{equation}
  \chi_{D}
  \qquad\text{and}\qquad
  \chi_{f}
\end{equation}
parameterize the strength of the suppression. In the modern RDCC framework
these amplitudes are no longer viewed as completely independent quantities.
Instead, they are expected to be controlled by the same relaxation dynamics
that generate the infrared parameter
\begin{equation}
  \alpha_{\mathrm{IR}}
  =
  0.0176 \pm 0.0021.
\end{equation}
A complete derivation of these amplitudes remains a subject for future work.

This modification introduces a scale dependence that reflects the relaxation
mechanism.

% ---------------------------------------------------------
\section{The RDCC Growth Index}
% ---------------------------------------------------------

The growth index is defined implicitly by
\begin{equation}
  f_{\mathrm{RDCC}}(k,a)
  =
  \Omega_{m}(a)^{\gamma_{\mathrm{RDCC}}(k,a)}.
\end{equation}
Solving for \(\gamma_{\mathrm{RDCC}}\) yields
\begin{equation}
  \gamma_{\mathrm{RDCC}}(k,a)
  =
  \frac{
    \ln f(a)
    +
    \ln
    \left[
      1
      -
      \chi_{f}
      \left(
        \frac{k}{k_{\mathrm{rel}}}
      \right)^{\alpha}
    \right]
  }{
    \ln \Omega_{m}(a)
  }.
\end{equation}
This expression shows that the growth index becomes scale-dependent and
acquires a distinctive RDCC signature near the relaxation scale.

% ---------------------------------------------------------
\section{Scale Dependence and RDCC Signatures}
% ---------------------------------------------------------

In the standard cosmological model the growth index is nearly constant. In
RDCC the scale dependence becomes
\begin{equation}
  \frac{\partial \gamma_{\mathrm{RDCC}}}{\partial k}
  \propto
  -
  \chi_{f}\,
  \alpha\,
  \frac{k^{\alpha-1}}{k_{\mathrm{rel}}^{\alpha}}.
\end{equation}
This modification produces
\begin{itemize}
  \item a suppressed growth index at small scales,
  \item an enhanced growth index at large scales,
  \item a characteristic turnover near \(k_{\mathrm{rel}}\).
\end{itemize}
These signatures provide a direct observational probe of the relaxation scale.

% ---------------------------------------------------------
\subsection{Connection to the Relaxation Scale}
% ---------------------------------------------------------

The growth hierarchy shares its characteristic scale with multiple
observational sectors of RDCC. The same relaxation scale appears in
\begin{itemize}
  \item tensor damping,
  \item stochastic gravitational-wave spectra,
  \item weak lensing,
  \item redshift-space distortions,
  \item gravitational slip.
\end{itemize}
The growth index therefore does not probe an isolated phenomenological
parameter. Instead it probes the same freeze-out geometry that governs the
global relaxation sector of the Universe.

% ---------------------------------------------------------
\section{Observational Probes}
% ---------------------------------------------------------

The growth index can be measured through several observables:
\begin{itemize}
  \item \textbf{Redshift-space distortions (RSD)} measure \(f\sigma_{8}\).
  \item \textbf{Weak lensing} measures the combination \(f\Omega_{m}\).
  \item \textbf{Velocity fields} measure the velocity divergence.
  \item \textbf{Cluster abundance} depends on the growth history.
\end{itemize}
In RDCC all of these observables acquire a scale-dependent modification that
reflects the relaxation mechanism. The growth index therefore provides a
dynamical fingerprint of the RDCC gravitational field.

% ---------------------------------------------------------
\section{Conclusion}
% ---------------------------------------------------------

The growth index of RDCC is reinterpreted as a consequence of the freeze-out
geometry of the relaxation sector. Rather than introducing the relaxation
scale as an independent phenomenological parameter, the present framework
adopts
\begin{equation}
  k_{\mathrm{rel}}
  =
  a_{\mathrm{fr}} H(a_{\mathrm{fr}}),
\end{equation}
where
\begin{equation}
  a_{\mathrm{fr}}
  \simeq
  10^{-4}
\end{equation}
denotes the freeze-out epoch.

The resulting causal chain is
\begin{equation}
  \alpha
  \;\rightarrow\;
  a_{\mathrm{fr}}
  \;\rightarrow\;
  k_{\mathrm{rel}}
  \;\rightarrow\;
  D(k,a)
  \;\rightarrow\;
  f(k,a)
  \;\rightarrow\;
  \gamma(k,a).
\end{equation}
The growth index therefore becomes a direct observable probe of the freeze-out
structure of the RDCC Universe.

Together with Companion XII, Companion VII, and Companion XLIV, the present
Companion establishes the connection between global relaxation dynamics and
the observable growth of cosmic structure.

\end{document}